\documentclass[10pt,twocolumn]{article}
\usepackage[margin=0.7in]{geometry}
\usepackage[T1]{fontenc}
\usepackage{lmodern}
\usepackage{microtype}
\usepackage{graphicx}
\usepackage{enumitem}
\usepackage{booktabs}
\usepackage{array}
\usepackage[dvipsnames]{xcolor}
\usepackage[hidelinks]{hyperref}
\setlist{nosep,leftmargin=*}
\setlength{\columnsep}{0.25in}
\setlength{\textfloatsep}{8pt}
\setlength{\abovecaptionskip}{3pt}
\setlength{\parskip}{2pt}
\makeatletter
\renewcommand\section{\@startsection{section}{1}{\z@}{6pt}{3pt}{\normalfont\large\bfseries}}
\renewcommand\subsection{\@startsection{subsection}{2}{\z@}{4pt}{2pt}{\normalfont\normalsize\bfseries}}
\renewcommand\paragraph{\@startsection{paragraph}{4}{\z@}{3pt}{-0.5em}{\normalfont\bfseries}}
\makeatother
\newcommand{\num}[1]{\textbf{#1}}
% An empty slot: names the experiment entry whose result will fill it.
\newcommand{\todo}[1]{\textcolor{BrickRed}{\textbf{[TODO: #1]}}}
% experiment entry id; hyphens may break
\newcommand{\entry}[1]{{\ttfamily\hyphenchar\font=`\-\relax #1}}
\setlength{\tabcolsep}{4pt}

\title{\vspace{-0.65in}\textbf{Repository Intelligence Layer: Final Report}\\[2pt]
\large Does the quality of a retriever's training data reach coding-agent success?}
\author{Sukhraj Singh\\ \small 2025202009 \and Neill Rawani\\ \small 2025202027 \and
Chaitanya Shah\\ \small 2025202033}
\date{\vspace{-6pt}IIIT Hyderabad}

\begin{document}
\maketitle
\vspace{-0.4in}

\todo{Abstract: written after the confirmatory contrasts and the agent runs are in; fills from
\entry{E-RQ2-FILTER}, \entry{E-RQ2-CTRL-SUBRAW}, \entry{E-RQ3-AGENT}.}

% =============================================================================================
\section{Introduction}

Coding agents work inside a fixed context window, while a real repository is many times larger.
Something has to decide which code the agent reads; when that choice is poor, the agent spends its
budget on irrelevant files or never sees the code it has to change. The \emph{Repository
Intelligence Layer} (RIL) is a repository-level retrieval layer: given an issue, it returns a small,
ranked set of Python code units that fit a token budget. The system is the apparatus; the
experiment is the study. Two things are built and only one is the contribution: the system is
engineering and is not novel on its own, while the experiment measures whether the way a
retriever's training pairs are handled changes whether the agent resolves the issue.

\paragraph{Research questions.}
\textbf{RQ1} Which retrieval signals contribute to repository-level code retrieval? Incremental
ablation on a dense baseline (dense, dense and BM25, dense and structure, all three), reported with
latency per configuration. This fixes the stack that the other questions hold constant.
\textbf{RQ2} How does training supervision quality and composition affect repository-level retrieval
under \emph{matched optimization budgets}? Four arms (\textsc{Raw}, \textsc{Filter}, \textsc{Weight},
\textsc{Curriculum}) trained for the same number of optimizer updates at matched batch size and
schedule; optimizer updates, distinct examples seen and effective epochs are reported for every cell,
and a random subsample of the raw pool of the filtered pool's size is trained beside every filtering
comparison. \textbf{RQ3} Does improving retriever supervision translate into improved coding-agent
performance? One fixed agent scaffold, one fixed base model, the same instances; the outcome is the
resolved rate, with, for every failure, whether the file of the reference fix was in the retrieved
context.

\paragraph{Claim.} Our literature audit did not identify a study that systematically evaluates the
controlled pipeline from retriever supervision quality, through retrieval performance, to
repository-level coding-agent success. The contribution is an empirical comparison, not an
intervention: filtering, weighting and curriculum each exist in the literature, and we compare them
under one protocol with the agent held fixed. The scope is Python only. Absolute resolved rates are
not statements about agent capability, because benchmark contamination raises them identically
across arms.

\paragraph{Contributions, stated as what is measured.} Every item points to a table below and to its
entry in the experiment registry; none is a claim about a retriever trained to convergence.
\begin{itemize}
  \item A function-level index with a typed code graph, and the published BM25 rows reproduced from
  it (Table~\ref{tab:bm25_repro}).
  \item Zero-shot dense encoders and rank fusion with BM25, and the graph's role measured against a
  same-file control (Tables~\ref{tab:dense_fusion}, \ref{tab:graph_role}).
  \item A noise dose-response of the supervision arms at matched updates, with the teacher caveat
  stated (Table~\ref{tab:noise_dose_response}).
  \item An unfiltered pull-request-mined training set, its rank profile against the published pool,
  and a first training pilot on it (Tables~\ref{tab:mined_vs_pool_ranks}, \ref{tab:mined_pilot}).
  \item The context assembly stage and how much of the reference fix lands in the context
  (Tables~\ref{tab:assembly_gold}, \ref{tab:assembly_cap}).
\end{itemize}
Every pilot result in this report is labelled \emph{exploratory} in its caption and is outside the
pre-registered contrasts of Section~\ref{sec:rq2}. \todo{Contributions on the confirmatory contrasts
and on resolved rate: \entry{E-RQ2-FILTER}, \entry{E-RQ3-AGENT}.}

% =============================================================================================
\section{Related work}

Only papers read in full are cited; each has a note in the project's reading log.

\paragraph{Retrievers for issue localization and the benchmarks.} SWE-bench defines the resolved-rate
task~\cite{swebench}. CORE-Bench provides repository-level retrieval labels at three levels and
a BM25 baseline~\cite{corebench}; SweRank and SweRank+ train code retrievers for issue
localization, with a consistency filter on the rank of the positive~\cite{swerank,swerankplus}; CoRNStack
supplies the contrastive training corpus that the published pool builds on~\cite{cornstack}. Two
follow-up benchmarks evaluate retrieval for coding agents and the cost of exploration
\cite{agentretrievalbench,costexploration}. One trained retrieval agent reports resolved rate
with the agent fixed~\cite{codegrep}.

\paragraph{Graph-structured localization.} Several systems give an agent or a retriever a repository
graph: LocAgent, KGCompass, SpIDER, GraphLocator and SemNav~\cite{locagent,kgcompass,spider,graphlocator,semnav};
RepoNavigator and LogicLoc make symbol resolution or program facts a tool~\cite{reponavigator,logicloc}.
The static call graphs underneath have measured precision and recall against dynamic
baselines for Python~\cite{pycg,dypybench}. We use the graph as a retrieval signal and test it
against a same-file control, which these systems do not report.

\paragraph{Training-data quality for retrievers.} Query filtering for code search~\cite{sun}; weighting
of negatives~\cite{rethinking}; synthetic data with a staged curriculum~\cite{coder}; false
negatives and pruning in text retrieval~\cite{rlhn}; a rank-relative-to-the-positive edit of mined
negatives~\cite{arhn}; loss-based detection of mismatched pairs~\cite{npc}; negative selection with
false-negative removal~\cite{nvretriever}; filtering the training pairs of a code retriever~\cite{map4cs};
and a study of which negative sources transfer~\cite{hardneg}. Curriculum and self-paced selection for
retrieval distillation~\cite{cldrd,penha,adam}. These works study one intervention each, in
most cases outside repository-level code retrieval; the comparison here is of the three
interventions under one matched protocol, followed through to the resolved rate.

\todo{Related-work refresh before submission: the dated search log and the papers still unread are
in the project's literature notes; nothing is cited from an abstract.}

% =============================================================================================
\section{System}
\label{sec:system}

\paragraph{Parser and typed graph.} The index is function level: every top-level function, method
and class is a unit, plus one residual module unit per file for module-level code; the unit schema is
frozen. Typed edges are \emph{contain}, \emph{import}, \emph{invoke}, \emph{inherit} and
\emph{override}; every edge carries a resolution status (resolved, ambiguous, external,
unresolved) and ambiguous call sites are dropped from the stored graph. Test and configuration edges
exist as exploratory families and are off by default. Rebuilds are byte-identical, and an
incremental updater re-reads only the files a commit changes and their dependants and equals the
full rebuild on every commit checked. \todo{Edge recovery against a dynamic oracle, and graph storage
cost: \entry{E-GRAPH-BUILD} tables.}

\paragraph{Retrieval signals.} BM25 in a reference form (stemmed tokens, a stop list, Lucene-variant
scoring) is the lexical signal. The dense signal is a code embedding model over a flat exact index
(no approximate search); the working encoder is CodeRankEmbed, and a larger encoder is used
zero-shot only because it cannot be fine-tuned on the reachable GPUs. The graph is a third signal in
the form of hop-decayed expansion from the top units.

\paragraph{Fusion.} Reciprocal rank fusion, which is rank-based, not a weighted score sum; the
fusion constant and the encoder cut were chosen on the development repositories.

\paragraph{Context assembly.} Candidates are taken in rank order and each unit's text is cut at a
per-unit cap; a unit that does not fit the remaining budget is skipped and the walk continues;
included blocks are ordered by file and line, each with a header giving path, lines, kind, name and
retrieval rank; the finished string is verified against the budget (ADR-014, provisional until the
agent pilot). Assembly is byte-identical across RQ2 arms, so that the only thing RQ3 varies is the
retriever. No graph or same-file expansion is applied in assembly. Table~\ref{tab:assembly_cap} shows
the cap measurement and Table~\ref{tab:assembly_gold} the gold in context. \todo{Budget in the agent
model's tokens: set by the agent pilot, \entry{E-RQ3-AGENT}.}

\paragraph{Latency harness.} Every commit that touches the query path appends a row of per-stage
latency (query encode, BM25, dense search, graph expansion, fusion, rerank, assemble, end to end) at
batch size one; latency is reported as p50, p95 and p99, never the mean, with warm-up calls
discarded, GPU sections synchronised and cold and warm cache separated. Indexing and training are
throughput, in separate tables, and numbers are never compared across hardware without saying so.
\todo{Per-stage latency table with the frozen configuration's benchmark snapshot, and nDCG per
millisecond: \entry{E-LAT-STAGES}, \entry{E-RQ1-ALL}.}

% =============================================================================================
\section{Evaluation setup}

\paragraph{Benchmark and levels.} CORE-Bench Level~2 asks for the file-and-function context needed to
edit; Level~3 asks for context beyond the edited files. Evaluation is at unit level (decision
D12): a chunk's positive counts for every unit it overlaps, the reading without module-residual
positives is primary and the benchmark-faithful reading is reported beside it (D18, provisional).
Chunk-level numbers with folder means are reported only where they reproduce the published rows
(Table~\ref{tab:bm25_repro}). The metrics are nDCG@10 at Level~2 and Recall@100 at Level~3.

\paragraph{Roles.} Every choice (K, teacher, weight form, update count, fusion constant, assembly
cap) is made on development repositories only: five at Level~2 and two to three at Level~3. SWE-bench
Verified is the final test and is looked at once, after the protocol, the cell list and the
decision rules are registered and the development analysis is complete. To raise the number of
repositories, which seeds cannot buy, eight further repositories take the test role under a rule
fixed before any result (decision D100): Python-class repositories outside the development set and
Verified with at least five Level~2 queries, ordered by a hash of the repository name, the first
eight taken, the next eight held as a reserve for a technical build failure only. They add Level~2
only; the Level~3 claim is stated on the fifteen repositories that have Level~3 labels (three
development, twelve Verified). \todo{Evaluation of the eight repositories and of Verified: one look,
after the development analysis; \entry{E-EVAL-REPOS-BUILD} built the stores, \entry{E-RQ2-FILTER} reads
them.}

\paragraph{Leakage.} The training pool shares no repository with CORE-Bench Level~2 and Level~3 or
Verified; instance-level disjointness is not enough. The held-out list is built first, current
repository names are resolved, the pool is filtered against it and the surviving counts are
reported. Tests assert repository disjointness, no patch or gold text in any training query,
that the quality-signal encoder is not the encoder being trained, and that no evaluation data is
reachable from a training path. \todo{Surviving counts for the mined set after the content-overlap
audit: \entry{E-MINE-FULL}.}

\paragraph{Statistics.} The primary number is the mean over queries. Intervals are 95\% bootstrap
intervals (2,000 draws, seed 0), clustered by repository where at least three repositories are
present and query-level otherwise (development Level~3 has two repositories, so its interval carries
query sampling only and is marked). Every contrast is a paired difference of two cells trained with
the same seed on the same queries; seed-to-seed spread is reported because the bootstrap carries
query and repository sampling, not retraining variance. Two cells that overlap within one seed
standard deviation are called indistinguishable at this sample size. Multiplicity: Holm step-down within
each pre-registered family, with the uncorrected intervals beside it. The smallest difference worth
reporting is $\delta = 1.5$ points at both levels; smaller differences are estimates with bounds,
never effects. The minimum detectable difference for $R$ repositories, $n$ queries and $k$ seeds is
$2.8016\sqrt{s^2_{\mathrm{repo}}/R + s^2_{\mathrm{query}}/n + (s^2_{\mathrm{seed\times query}}/n +
s^2_{\mathrm{seed}})/k}$ with variance components estimated from stored runs; seeds divide only the
last bracket, so repositories, not seeds, reduce the repository term.

\paragraph{Evidence classes.} Every number carries one: \emph{measured} (a run in this project, with
log, configuration hash, seed and hardware), \emph{operator\_reported}, \emph{reasoned} (arithmetic
without new data) or \emph{synthetic}; numbers quoted from a paper are \emph{paper\_reported}.

% =============================================================================================
\section{RQ1: which signals matter}

Table~\ref{tab:dense_fusion} gives the zero-shot signals on the final-test-style reading, and
Table~\ref{tab:graph_role} the graph against the same-file control: expansion over typed edges is
always compared with taking more units of the files already found, at the same candidate budget,
because a plain expansion baseline understates what file-level information already provides. The
tables are exploratory, with settings chosen on the development repositories.

\begin{table*}[t]
\centering\small
\caption{BM25 reproduction of the published baseline rows at chunk level, mean over folders (the reading that reproduces them). Evidence class: \emph{measured; the paper columns are paper\_reported}. n is the folder and query count of each row; the query-level mean over the same queries is in the source table and differs. Entry \entry{E-BM25-REPRO}.}
\label{tab:bm25_repro}
\begin{tabular}{lrrrrrr}
\toprule
Set & Folders & Queries & nDCG@10 & nDCG@10 paper & Recall@100 & Recall@100 paper \\
\midrule
Verified L2 & 12 & 432 & 11.94 & 11.9 & 46.74 & 46.7 \\
Verified L3 & 12 & 302 & 22.19 & 22.2 & 33.89 & 33.9 \\
All sub-datasets L2 & 632 & 5061 & 13.00 & 13.0 & 38.82 & 38.8 \\
\bottomrule
\end{tabular}
\end{table*}

\begin{table*}[t]
\centering\small
\caption{Zero-shot dense encoders and rank fusion with BM25 (final-test-style reading). Evidence class: \emph{measured, exploratory}. Verified-six, 83 L2 and 55 L3 queries in 6 repositories; the interval is repository-clustered; cut and k were chosen on the development repositories. Entry \entry{E-DENSE-FUSION}.}
\label{tab:dense_fusion}
\begin{tabular}{lrrrr}
\toprule
System & Cut & RRF k & L2 nDCG@10 & L3 Recall@100 \\
\midrule
BM25 &  &  & 18.0 [12.1, 25.4] & 31.8 [23.8, 51.0] \\
CodeRankEmbed & 1024 &  & 17.9 [13.6, 21.2] & 33.6 [23.6, 51.0] \\
BM25 + CodeRankEmbed (RRF) & 1024 & 100 & 23.0 [13.4, 31.8] & 38.0 [27.7, 58.4] \\
Qwen3-Embedding-0.6B & 512 &  & 17.0 [12.2, 24.2] & 39.1 [31.6, 54.1] \\
BM25 + Qwen3-Embedding-0.6B (RRF) & 512 & 20 & 24.9 [18.4, 30.9] & 41.5 [33.1, 59.8] \\
\bottomrule
\end{tabular}
\end{table*}

\begin{table*}[t]
\centering\small
\caption{Graph expansion against the same-file control on a fused first stage, equal candidate budget. Evidence class: \emph{measured, exploratory}. Verified-six: 83 L2 and 55 L3 queries in 6 repositories; dev: 351 and 155 queries in 5 and 2 repositories. Intervals are repository-clustered, except $^q$ (two repositories: query-level only). Entry \entry{E-GRAPH-ROLE}.}
\label{tab:graph_role}
\begin{tabular}{>{\raggedright\arraybackslash}p{0.28\linewidth}>{\raggedleft\arraybackslash}p{0.16\linewidth}>{\raggedleft\arraybackslash}p{0.16\linewidth}>{\raggedleft\arraybackslash}p{0.16\linewidth}>{\raggedleft\arraybackslash}p{0.16\linewidth}}
\toprule
Candidate source & Verified-six L2 nDCG@10 & Verified-six L3 Recall@100 & dev L2 nDCG@10 & dev L3 Recall@100 \\
\midrule
first stage, RRF of BM25 and CodeRankEmbed (value) & 23.0 & 37.8 & 25.1 & 42.1 \\
\midrule
graph as stored: change against first stage & +1.9 [0.1, 4.8] & +5.7 [1.3, 7.8] & +1.9 [0.8, 3.7] & +12.2 [10.1, 14.4]$^q$ \\
same file, by first-stage score (control): change against first stage & +0.5 [$-$0.6, 2.5] & +7.7 [6.5, 8.3] & +1.2 [$-$0.3, 3.2] & +13.2 [11.0, 15.6]$^q$ \\
same file, random order: change against first stage & $-$1.7 [$-$3.0, 0.2] & +2.5 [0.9, 5.3] & +0.9 [$-$1.0, 3.2] & +11.3 [9.1, 13.7]$^q$ \\
next best units, no file constraint: change against first stage & $-$0.4 [$-$1.4, 0.5] & $-$0.2 [$-$0.9, 0.2] & $-$0.6 [$-$1.5, 0.6] & +0.5 [0.1, 0.9]$^q$ \\
graph, non-contain edges only: change against first stage & +1.5 [$-$0.1, 5.2] & +0.4 [$-$1.7, 1.3] & +0.4 [$-$0.4, 1.6] & +5.7 [4.0, 7.5]$^q$ \\
graph, contain edges only: change against first stage & +1.6 [0.3, 2.7] & +5.0 [3.0, 8.3] & +2.2 [0.8, 4.2] & +13.1 [10.6, 15.5]$^q$ \\
\midrule
graph as stored: minus the same-file control & +1.4 [0.2, 3.1] & $-$2.0 [$-$6.3, 0.6] & +0.7 [0.1, 1.3] & $-$1.0 [$-$3.3, 1.2]$^q$ \\
graph, non-contain edges only: minus the same-file control & +0.8 [$-$1.0, 3.6] & $-$7.3 [$-$9.3, $-$6.6] & $-$0.9 [$-$1.9, $-$0.2] & $-$6.1 [$-$8.5, $-$3.9]$^q$ \\
graph, contain edges only: minus the same-file control & +1.2 [0.6, 1.4] & $-$2.0 [$-$4.3, 2.9] & +0.4 [$-$0.7, 1.2] & +0.9 [$-$0.9, 2.6]$^q$ \\
\bottomrule
\end{tabular}
\end{table*}


\paragraph{Reading.} \todo{One paragraph per table once the RQ1 ladder exists: the contribution of
BM25 over dense, and of structure over the same-file control, at both levels, with the interval
and the latency column; \entry{E-RQ1-DENSE}, \entry{E-RQ1-DENSE-BM25}, \entry{E-RQ1-DENSE-STRUCT},
\entry{E-RQ1-ALL}.}

\paragraph{Open graph question.} Typed edges that leave a file are tested against the control in
Table~\ref{tab:graph_role}; whether training-time structure or an agent-tool role adds anything is
unmeasured. \todo{Stacked expansion at the same budget on development, Verified-six and the third
development Level~3 repository: \entry{E-GRAPH-L3}.}

% =============================================================================================
\section{RQ2: supervision under matched budgets}
\label{sec:rq2}

\subsection{Protocol}

The retrieval stack is frozen across arms; only the supervision strategy varies. Held identical:
base checkpoint (CodeRankEmbed), architecture, pooling, batch shape, learning-rate schedule,
temperature, maximum length, one negative policy for all arms, the evaluation set, the flat exact
index, rank fusion and byte-identical context assembly. The protocol settles fifteen items before any
final run (O1 to O15); the ones that shape the claim are:
\begin{itemize}
  \item \textbf{Updates (O1):} the primary budget is 100 optimizer updates of a longer schedule, with
  1,000 updates as a sensitivity check on a subset of cells, because the pilots saturate early.
  \item \textbf{Signal and teachers (O3):} the quality signal is SweRank's K-rank (minimum rank of
  the positive under a base retriever) at K = 20, positive level. BM25 is the primary teacher because the
  published pool's construction does not favour it; SweRankEmbed-Small is reported with the caveat
  that it shaped the pool. No cell uses the encoder being trained as its own teacher.
  \item \textbf{Confirmatory scope (O2):} exactly two contrasts, both on the unfiltered mined set:
  \textsc{Filter} (BM25 teacher) against \textsc{Raw}, and against the size-matched random
  subsample. Everything else (the second teacher, \textsc{Weight}, \textsc{Curriculum}, the injected-noise
  calibration, the dose-response, the noise-audit prediction) is exploratory and labelled so.
  \item \textbf{Matching rule:} equal optimizer updates and reported distinct-example counts; a
  cell that cannot reach the common count without repeating examples is flagged, never silently
  included.
  \item \textbf{Design of record (O7, O12):} three seeds per cell and as many evaluation
  repositories as the budget allows (the eight above), over more seeds.
  \item \textbf{Fresh noise (O15):} confirmatory runs draw new noise realisations with the same
  rule; the pilot realisations are not reused.
  \item \textbf{Retention (O9):} final models of runs that feed RQ3 are retained (half precision).
\end{itemize}

\paragraph{Decision rules, fixed before any final run.} \emph{Positive}: the lower end of the
adjusted interval is above zero at Level~2 (repository-clustered) and at Level~3 on the final test,
with the same sign on development at both levels, and the same holds against the subsample for the claim
``quality beyond quantity''. \emph{Negative, with a bound}: the interval includes zero and its upper
end is below $\delta$ at both levels, reported as ``no gain larger than $u$ points at this
design''. \emph{Inconclusive}: the lower end is at most zero and the upper end is at least $\delta$
at either level, reported with the minimum detectable difference. For the injected-noise cells,
\textsc{Raw} must degrade against the clean cell and \textsc{Oracle-Clean} must recover it, or the
calibration is reported as not informative. A null stays a null: if arms differ in retrieval but
not in resolved rate, RQ3 reports the retrieval-versus-reasoning failure breakdown.

\subsection{Pilots: dose-response and the mined set}

Table~\ref{tab:noise_dose_response} injects label noise of a known kind at a known rate into the
published pool and measures what the arms recover; Table~\ref{tab:mined_vs_pool_ranks} compares the
rank profile of the unfiltered mined set with the published pool's; Table~\ref{tab:mined_pilot} is the
first training run on the mined set. Two cautions are part of the result. The published pool looks
pre-selected by one retriever, so \textsc{Raw} on it is not raw; and injected noise is a proxy whose
transfer to real noise is open (a second kind of noise failed its pre-registered sanity check).
\todo{Sibling-noise decision and its two further seeds: \entry{E-PILOT-NOISE-SIB}.}

\begin{table*}[t]
\centering\small
\caption{Injected-noise dose-response: supervision arms against RAW at matched updates. Evidence class: \emph{measured, exploratory}. Development repositories: 351 L2 queries in 5 repositories, 155 L3 queries in 2; L2 intervals are repository-clustered, L3 intervals are query-level (two repositories). Share of oracle is the arm's change divided by the ORACLE-CLEAN change. Injected noise is a proxy for real label noise. $^\dagger$SweRankEmbed-Small shaped the selection of the public pool (E-POOL-RANKS), so its cells are the favoured teacher; the BM25 cells are the ones the confirmatory protocol rests on. Entry \entry{E-PILOT-NOISE}.}
\label{tab:noise_dose_response}
\begin{tabular}{>{\raggedright\arraybackslash}p{0.3\linewidth}>{\raggedleft\arraybackslash}p{0.06\linewidth}>{\raggedleft\arraybackslash}p{0.17\linewidth}>{\raggedleft\arraybackslash}p{0.09\linewidth}>{\raggedleft\arraybackslash}p{0.17\linewidth}>{\raggedleft\arraybackslash}p{0.09\linewidth}}
\toprule
Arm minus RAW at the same r & Seeds & L2 nDCG@10 change & Share of oracle & L3 Recall@100 change & Share of oracle \\
\midrule
r = 0.2 &  &  &  &  &  \\
RAW minus clean RAW & 2 & $-$1.7 [$-$2.4, $-$0.9] &  & $-$1.1 [$-$1.6, $-$0.6] &  \\
ORACLE-CLEAN & 2 & +0.9 [0.4, 1.6] &  & +2.2 [1.4, 3.0] &  \\
FILTER, BM25 teacher & 2 & +1.0 [$-$0.7, 2.9] & 1.09 & +2.1 [1.1, 3.1] & 0.93 \\
FILTER, SweRankEmbed-Small teacher$^\dagger$ & 2 & +1.6 [0.7, 2.6] & 1.68 & +1.7 [0.9, 2.5] & 0.75 \\
WEIGHT, SweRankEmbed-Small teacher$^\dagger$ & 2 & +0.8 [0.3, 1.2] & 0.85 & +0.7 [0.3, 1.0] & 0.31 \\
size-matched subsample & 2 & $-$0.2 [$-$0.8, 0.6] & $-$0.20 & +0.6 [0.1, 1.2] & 0.29 \\
r = 0.4 &  &  &  &  &  \\
RAW minus clean RAW & 3 & $-$3.5 [$-$4.9, $-$1.6] &  & $-$1.9 [$-$2.7, $-$1.0] &  \\
ORACLE-CLEAN & 3 & +2.7 [1.3, 3.9] &  & +3.0 [1.9, 4.0] &  \\
FILTER, BM25 teacher & 3 & +1.8 [0.4, 2.9] & 0.68 & +1.2 [0.2, 2.1] & 0.40 \\
FILTER, SweRankEmbed-Small teacher$^\dagger$ & 3 & +3.0 [1.1, 4.3] & 1.13 & +2.6 [1.6, 3.6] & 0.87 \\
WEIGHT, SweRankEmbed-Small teacher$^\dagger$ & 3 & +1.9 [0.8, 2.7] & 0.70 & +1.1 [0.6, 1.7] & 0.38 \\
size-matched subsample & 3 & $-$0.2 [$-$0.6, 0.0] & $-$0.09 & $-$1.0 [$-$2.2, $-$0.1] & $-$0.34 \\
\midrule
minimum detectable difference, r = 0.4, with the repository term & 3 & 2.16 &  & 4.88 &  \\
minimum detectable difference, r = 0.4, query level & 3 & 1.40 &  & 2.20 &  \\
\bottomrule
\end{tabular}
\end{table*}

\begin{table*}[t]
\centering\small
\caption{Share of directories without a positive in the teacher's top K: pull-request-mined set against the published pool. Evidence class: \emph{measured, exploratory}. Pool-like reading of the mined pilot set: 300 directories with a positive (fewer where a teacher did not finish, see the Mined column); Wilson intervals; the pool column is the pool estimate from a stratified sample of directories; one seed, one pilot. Entry \entry{E-MINE-PILOT}, \entry{E-POOL-RANKS}.}
\label{tab:mined_vs_pool_ranks}
\begin{tabular}{lrrrr}
\toprule
Teacher & K & Mined directories & Mined set, no positive in top K (\%) & Published pool (\%) \\
\midrule
BM25 & 10 & 300 & 53.3 [47.7, 58.9] & 23.2 [16.4, 31.0] \\
BM25 & 20 & 300 & 45.3 [39.8, 51.0] & 15.1 [9.4, 21.1] \\
BM25 & 40 & 300 & 36.3 [31.1, 41.9] & 13.4 [8.0, 19.6] \\
CodeRankEmbed & 10 & 254 & 40.2 [34.3, 46.3] & 15.3 [8.8, 22.2] \\
CodeRankEmbed & 20 & 254 & 32.3 [26.8, 38.3] & 9.9 [5.2, 15.9] \\
CodeRankEmbed & 40 & 254 & 25.2 [20.2, 30.9] & 4.3 [1.6, 8.0] \\
SweRankEmbed-Small & 10 & 300 & 31.3 [26.4, 36.8] & 1.7 [0.2, 3.9] \\
SweRankEmbed-Small & 20 & 300 & 26.3 [21.7, 31.6] & 1.2 [0.0, 3.5] \\
SweRankEmbed-Small & 40 & 300 & 21.3 [17.1, 26.3] & 1.2 [0.0, 3.5] \\
\bottomrule
\end{tabular}
\end{table*}

\begin{table*}[t]
\centering\small
\caption{Training on the pull-request-mined set: gains over the untrained encoder and arms. Evidence class: \emph{measured, exploratory}. Development repositories: 351 L2 queries in 5 repositories, 155 L3 queries in 2; L2 intervals are repository-clustered, L3 query-level; mined set and pool both at a matched number of optimizer updates and distinct examples. Intervals carry query and repository sampling, not retraining variance; at this number of seeds a difference is read against the minimum detectable difference in the last row. Entry \entry{E-MINED-PILOT}.}
\label{tab:mined_pilot}
\begin{tabular}{>{\raggedright\arraybackslash}p{0.4\linewidth}>{\raggedleft\arraybackslash}p{0.07\linewidth}>{\raggedleft\arraybackslash}p{0.2\linewidth}>{\raggedleft\arraybackslash}p{0.2\linewidth}}
\toprule
Contrast & Seeds & L2 nDCG@10 & L3 Recall@100 \\
\midrule
RAW on the mined set, minus the untrained encoder & 2 & +7.9 [5.0, 10.1] & +12.2 [9.9, 14.4] \\
RAW on the published pool, minus the untrained encoder & 2 & +9.6 [6.9, 11.6] & +11.1 [9.1, 13.2] \\
RAW on the mined set minus RAW on the pool & 2 & $-$1.7 [$-$3.1, $-$0.2] & +1.1 [$-$0.2, 2.3] \\
FILTER (BM25 teacher) minus RAW, mined set & 2 & +0.9 [0.5, 1.3] & 0.0 [$-$0.9, 0.8] \\
size-matched subsample minus RAW, mined set & 2 & $-$0.2 [$-$0.7, 0.4] & $-$0.2 [$-$0.9, 0.5] \\
\midrule
minimum detectable difference, FILTER minus RAW, with repository term & 2 & 1.44 & 2.36 \\
\bottomrule
\end{tabular}
\end{table*}


\subsection{Noise audit}

The share of mined pairs that carry harmful noise, and of what kind, is estimated from a
stratified sample: a wrong positive (an edited function the issue is not about), a sibling of the
relevant function, or a missing positive (a false negative among the negatives, which is not
harmful on the positive side). The sample is stratified by K-rank (kept and dropped under each
teacher) so that the precision and recall of the filter are estimated, not only the overall rate. Two
readers who are not the lead read the pairs, so the lead's knowledge of the arms cannot influence the
sample; a 100-pair overlap gives their agreement. A binomial rate to a half-width of 0.05 needs
about 400 pairs. If the audit cannot be done the report states the K-rank profile of the mined set
as a proxy and does not test the prediction from the audited rate. \todo{Audited noise rate and
reader agreement: \entry{E-LABEL-AUDIT-SHEET}.}

\subsection{Confirmatory contrasts}

\todo{\textsc{Filter} (BM25) minus \textsc{Raw} on the mined set, both levels, three seeds, adjusted
and unadjusted intervals, verdict by the rules above: \entry{E-RQ2-FILTER}, \entry{E-RQ2-RAW}.}
\todo{\textsc{Filter} (BM25) minus the size-matched subsample, same layout, with the distinct-example
counts of every cell: \entry{E-RQ2-CTRL-SUBRAW}.}
\todo{Exploratory arms (second teacher, \textsc{Weight}, \textsc{Curriculum}) and the dose-response
at fresh draws: \entry{E-RQ2-WEIGHT}, \entry{E-RQ2-CURRICULUM}.}
\todo{Mined-set size after the full mining run, yield and surviving counts: \entry{E-MINE-FULL}.}
\todo{Longer-schedule sensitivity cells and the quantity--quality curve if the core completes:
\entry{E-RQ2-SEC-QQ}.}

% =============================================================================================
\section{RQ3: from retrieval to resolved rate}

\paragraph{Design.} Each RQ2 retriever is plugged into one fixed agent scaffold with one fixed base
model, temperature and seed, over the same instances; the outcome is the resolved rate with a paired
McNemar test and intervals, and for every failure the gold-in-context indicator that separates
retrieval failures from reasoning failures. The scaffold and the base model are open decisions
(ADR-007, ADR-008) and nothing is run until their cost and noise are known. The assembly policy
(Section~\ref{sec:system}, Tables~\ref{tab:assembly_gold}, \ref{tab:assembly_cap}) is shared by all arms.

\begin{table*}[t]
\centering\small
\caption{Gold in the assembled context by token budget and first stage. Evidence class: \emph{measured, exploratory}. Verified-six, 83 L2 and 55 L3 queries in 6 repositories; budget in approximate tokens (not the agent model's tokenizer); repository-clustered intervals; unit cap as in the cap table; read once. Entry \entry{E-ASSEMBLY}.}
\label{tab:assembly_gold}
\begin{tabular}{>{\raggedright\arraybackslash}p{0.19\linewidth}>{\raggedleft\arraybackslash}p{0.07\linewidth}>{\raggedleft\arraybackslash}p{0.15\linewidth}>{\raggedleft\arraybackslash}p{0.15\linewidth}>{\raggedleft\arraybackslash}p{0.15\linewidth}>{\raggedleft\arraybackslash}p{0.15\linewidth}}
\toprule
First stage & Budget & L2 gold files in context (\%) & L2 queries with no gold file (\%) & L3 gold files in context (\%) & L3 queries with no gold file (\%) \\
\midrule
BM25 & 4000 & 51.2 [41.4, 69.4] & 45.8 [28.6, 56.3] & 46.4 [38.5, 64.4] & 21.8 [3.3, 33.9] \\
BM25 & 8000 & 64.3 [55.0, 86.7] & 31.3 [12.2, 39.7] & 59.7 [54.2, 77.6] & 9.1 [0.0, 14.3] \\
BM25 & 16000 & 69.7 [62.4, 87.5] & 26.5 [10.3, 33.8] & 65.8 [60.0, 83.0] & 5.5 [0.0, 7.9] \\
BM25 & 32000 & 81.5 [72.1, 95.9] & 16.9 [4.1, 26.7] & 76.2 [69.6, 91.6] & 3.6 [0.0, 7.1] \\
\midrule
RRF of BM25 and CodeRankEmbed & 4000 & 64.3 [52.5, 82.1] & 31.3 [16.7, 43.4] & 57.7 [47.0, 73.9] & 14.6 [0.0, 25.0] \\
RRF of BM25 and CodeRankEmbed & 8000 & 74.5 [63.5, 92.9] & 21.7 [6.1, 32.5] & 70.0 [60.2, 86.8] & 3.6 [0.0, 7.1] \\
RRF of BM25 and CodeRankEmbed & 16000 & 82.1 [73.0, 98.0] & 14.5 [2.0, 22.4] & 76.9 [69.8, 90.8] & 0.0 [0.0, 0.0] \\
RRF of BM25 and CodeRankEmbed & 32000 & 90.0 [83.9, 100.0] & 8.4 [0.0, 14.0] & 84.4 [78.6, 94.0] & 0.0 [0.0, 0.0] \\
\bottomrule
\end{tabular}
\end{table*}

\begin{table}[t]
\centering\small
\caption{Per-unit cap of the assembled context: share of gold units kept intact. Evidence class: \emph{measured, exploratory}. Development repositories only, mean over 8 cells of level, budget and first stage (351 L2 queries in 5 repositories in an L2 cell); the rule was fixed before the run. Entry \entry{E-ASSEMBLY-CAP}.}
\label{tab:assembly_cap}
\begin{tabular}{>{\raggedleft\arraybackslash}p{0.2\linewidth}>{\raggedleft\arraybackslash}p{0.2\linewidth}>{\raggedleft\arraybackslash}p{0.2\linewidth}>{\raggedright\arraybackslash}p{0.16\linewidth}}
\toprule
Unit cap (approximate tokens) & Gold units intact (\%) & Gold files in context (\%) & Selected \\
\midrule
500 & 25.61 & 76.60 & no \\
1000 & 28.70 & 74.95 & yes \\
2000 & 29.35 & 74.12 & no \\
\bottomrule
\end{tabular}
\end{table}


\paragraph{Power.} With the published fixed-agent retrieval swap read in full as the only comparison
available, the differences it reports are below the minimum detectable effect at Verified sizes, so
the resolved-rate comparison may be reportable only as an estimate with an interval. The smallest
worthwhile difference is stated before the runs. \todo{Power table from stored discordance and
run-to-run spread of the agent pilot: \entry{E-POWER}.}

\paragraph{Fallback (declared before any final run).} If the arms do not separate under the rule
above, RQ3 contrasts \textsc{Raw} against the untrained encoder and is labelled as measuring
``does training help'', not ``does supervision quality help''. \todo{Agent runs, resolved rate with
interval, McNemar, failure breakdown and resolved rate per agent wall-clock: \entry{E-RQ3-AGENT}.}

% =============================================================================================
\section{Threats to validity}
\label{sec:threats}

\begin{table*}[t]
\centering\small
\caption{Threats to validity and their mitigations. Evidence class: \emph{reasoned}, with the
evidence behind each row in the entries named.}
\label{tab:threats}
\begin{tabular}{@{}p{0.04\textwidth}p{0.35\textwidth}p{0.53\textwidth}@{}}
\toprule
& Threat & What we do about it \\
\midrule
T1 & The published pool was built with a retriever in the loop, so \textsc{Raw} on it is not raw
(\entry{E-POOL-RANKS}) & Unfiltered mined set as the pool on which \textsc{Raw} is defined; the pool is kept as the
comparison population so the selection stays measurable \\
T2 & Pretraining exposure: the favoured teacher was trained on the same repository family; the base
encoder and the agent's model may have seen the repositories (\entry{E-CONTAM-SCAN}, \entry{E-CONTAM-EDITED}) &
BM25 carries the primary claim; each teacher reported separately; arms compared inside one benchmark, so inflation is equal across arms \\
T3 & Injected noise is a proxy; a neighbouring-function kind failed its sanity check
(\entry{E-PILOT-NOISE-SIB}) & Noise audit of the mined set; ``findable'' is measured, not assumed; the dose-response is a calibration, not a model of real labels \\
T4 & No trained checkpoint of a pilot result run is retained, and RQ3 needs the models & Retention switch in
the training jobs; final models of runs that feed RQ3 archived (\entry{E-RETAIN-FINAL}) \\
T5 & Few repositories: five at Level~2 and two to three at Level~3 on development; one Verified project dominates its Level~3 labels & Repository-clustered intervals
where possible, the repository term beside every query-level interval, a leave-one-project-out sensitivity, eight added Level~2 repositories \\
T6 & Short training: the pilots run a small fraction of the master schedule and could miss an effect that appears later & Longer-schedule sensitivity cells on the mined set (needs a mined pool large enough to avoid repeats) \\
T7 & Data-budget confound: filtering sees fewer distinct examples at equal updates & The size-matched subsample is in every comparison, one draw per seed, distinct-example counts reported \\
T8 & A teacher equal to the encoder being trained is a self-confirming loop & Removed from the final cells; a separately pretrained teacher only \\
T9 & Multiplicity and analytic flexibility: many contrasts were run in the pilots without correction & One confirmatory family fixed in advance, Holm within families, rules written before any final run \\
T10 & Single language, single encoder; Level~3 labels are judged by models & Claims scoped to Python and to the CodeRankEmbed base; Level~3 is secondary evidence for the context question \\
T11 & Level~3 rests on fifteen repositories, twelve of which are the final test & The Level~3 statement names its repositories; Level~2 carries the repository-count argument \\
\bottomrule
\end{tabular}
\end{table*}

Table~\ref{tab:threats} lists the threats that the protocol raises itself. Two more are standing
limits of the evidence: absolute levels are not capability claims, and every retrieval result here
is for Python.
\todo{Threats that the agent runs add (scaffold sensitivity, model non-determinism, repository
contamination of the agent's model): \entry{E-RQ3-AGENT}.}

% =============================================================================================
\section{Conclusion}

What is measured so far is the apparatus and the pilots: the published BM25 rows are reproduced,
the zero-shot signals and fusion are measured, the graph is compared with a same-file control, the
supervision arms are compared under injected noise at matched updates, and the mined set and the
assembly stage are built and measured once (Tables~\ref{tab:bm25_repro}--\ref{tab:assembly_cap}). Each
of these is exploratory, on development repositories unless the caption says Verified-six, and none is a
confirmatory result. The confirmatory question, whether \textsc{Filter} with the BM25 teacher
beats \textsc{Raw} and the size-matched subsample on unfiltered pairs at matched budgets, and the
resolved-rate question, are not answered here. \todo{Conclusion of the confirmatory contrasts, stated
with its bound and its verdict, null or not: \entry{E-RQ2-FILTER}, \entry{E-RQ3-AGENT}.}

% =============================================================================================
\begin{thebibliography}{99}\small
\bibitem{swebench} % note: swe-bench-2310.06770
C.~Jimenez, J.~Yang et al. SWE-bench: can language models resolve real-world GitHub issues? ICLR 2024 (arXiv 2310.06770).
\bibitem{corebench} % note: core-bench-2606.11864
Zhang et al. CORE-Bench (arXiv 2606.11864, 2026).
\bibitem{swerank} % note: swerank-2505.07849
Gangi Reddy, Suresh, Doo et al. SweRank (arXiv 2505.07849, ICLR 2026 version).
\bibitem{swerankplus} % note: swerank-plus-2512.20482
Gangi Reddy, Liu et al. SweRank+ (arXiv 2512.20482).
\bibitem{cornstack} % note: cornstack-2412.01007
Suresh, Gangi Reddy et al. CoRNStack (arXiv 2412.01007, ICLR 2025).
\bibitem{agentretrievalbench} % note: agent-retrieval-bench-2607.24882
Qin and Xie. Agent Retrieval Bench: evaluating repository context retrieval for coding agents (arXiv 2607.24882, 2026).
\bibitem{costexploration} % note: cost-effective-exploration-2608.29675
Al Awad and Ivanov. Cost-effective repository exploration for agentic issue localization (arXiv 2608.29675, 2026).
\bibitem{codegrep} % note: codegrep-2608.05886
Chen et al. CodeGrep: an RL-trained retrieval agent for coding agents (arXiv 2608.05886, 2026).
\bibitem{locagent} % note: locagent-2503.09089
Chen, Tang, Deng et al. LocAgent (arXiv 2503.09089, ACL 2025).
\bibitem{kgcompass} % note: kgcompass-2503.21710
Yang, Ren, Jin et al. KGCompass (arXiv 2503.21710).
\bibitem{spider} % note: spider-2512.16956
Chaudhari et al. SpIDER (arXiv 2512.16956).
\bibitem{graphlocator} % note: graphlocator-2512.22469
Liu et al. GraphLocator: graph-guided causal reasoning for issue localization (arXiv 2512.22469).
\bibitem{semnav} % note: semnav-2609.31176
Wei, Lei and Li. SemNav: semantic navigation for repository-level issue localization (arXiv 2609.31176, 2026).
\bibitem{reponavigator} % note: reponavigator-2512.20957
Zhang, Duan, Zhang et al. One tool is enough: reinforcement learning of agents for repository-level code navigation (arXiv 2512.20957).
\bibitem{logicloc} % note: logicloc-2604.16021
Xu, Wu, Zhang and Li. Neurosymbolic repo-level code localization (arXiv 2604.16021, 2026).
\bibitem{pycg} % note: pycg-2103.00587
Salis et al. PyCG: practical call graph generation in Python. ICSE 2021 (arXiv 2103.00587).
\bibitem{dypybench} % note: dypybench-2403.00539
Bouzenia, Krishan and Pradel. DyPyBench. FSE 2024 (arXiv 2403.00539).
\bibitem{sun} % note: sun-2202.06649
Sun et al. On the importance of building high-quality training datasets for neural code search. ICSE 2022 (arXiv 2202.06649).
\bibitem{rethinking} % note: rethinking-negative-pairs-2310.08069
Li et al. Rethinking negative pairs in code search. EMNLP 2023 (arXiv 2310.08069).
\bibitem{coder} % note: coder-2505.12697
Li et al. Towards a generalist code embedding model based on massive data synthesis (arXiv 2505.12697).
\bibitem{rlhn} % note: rlhn-2505.16967
Thakur, Zhang, Ma and Lin. Hard negatives, hard lessons: revisiting training data quality for robust information retrieval. EMNLP 2025 Findings (arXiv 2505.16967).
\bibitem{arhn} % note: arhn-2604.11092
Choi et al. Answer-centric relabeling of hard negatives with open-source models for dense retrieval (arXiv 2604.11092, 2026).
\bibitem{npc} % note: npc-2311.03798
Zhang et al. Noisy pair corrector for dense retrieval. Findings of EMNLP 2023 (arXiv 2311.03798).
\bibitem{nvretriever} % note: nv-retriever-2407.15831
Moreira et al. NV-Retriever: improving text embedding models with effective hard-negative mining (arXiv 2407.15831).
\bibitem{map4cs} % note: map4cs-2610.11727
Chen et al. MAP4CS: a multi-dimensional data pruning framework for efficient code retriever fine-tuning (arXiv 2610.11727, 2026).
\bibitem{hardneg} % note: hard-negative-transfer-2603.20990
Sinha, Seetharaman and Bansal. What tracks hard-negative transfer? A study of lexical and semantic signals (arXiv 2603.20990).
\bibitem{cldrd} % note: cl-drd-2204.13679
Zeng, Zamani and Vinay. Curriculum learning for dense retrieval distillation. SIGIR 2022 (arXiv 2204.13679).
\bibitem{penha} % note: curriculum-ir-1912.08555
Penha and Hauff. Curriculum learning strategies for IR. ECIR 2020 (arXiv 1912.08555).
\bibitem{adam} % note: adam-2212.10192
Tao et al. Dense retrieval distillation with adaptive dark examples (arXiv 2212.10192).
\end{thebibliography}

\end{document}
